\chapter*{Preface}
Compute $1^3+2^3+\dots+10^3$. The sum is $3025$, and so is $(1+2+\dots+10)^2$.
Ten more values behave the same way. Is the pattern a coincidence, or a fact about
every positive integer? Further checks still only check further cases. A proof has
to say why the pattern cannot fail, and it has to say so in a way a reader can test.

That is the skill this volume teaches. \emph{Order and Chance} follows number
theory, algebra and probability from their first principles toward the questions
researchers ask where these subjects meet. Mathematical maturity is the ability
to understand a statement, find an argument, calculate with a purpose, and decide
what has actually been proved.

You need secondary-school algebra: arithmetic with fractions, powers, equations
and inequalities, and familiarity with graphs and elementary geometry. You do not
need a previous course in proof, calculus, linear algebra or probability. The
language chapter starts with what a mathematical assertion means; the later
chapters rebuild their tools before using them. Some three-star problems will
require sustained thought. Difficulty is part of the work, not evidence that you
have missed a prerequisite.

\section*{The route through the volume}
Each chapter opens with one concrete question and then builds the tool that
answers it. The questions, in order, are: what does a mathematical sentence
promise; why does a pattern in the first twenty integers hold for all of them;
how do overlapping counts stay exact; what does it mean for infinitely many
numbers to have a sum; how does multiplication turn the plane; when does a
polynomial identity follow from finitely many checks; what does a matrix know
that its entries hide; how large is the error in replacing a function by a
polynomial; which estimates are forced by convexity; and what does a finite
experiment say about the long run.

There are ten chapters, numbered from zero. Chapter~\ref{ch:language} introduces
statements, quantifiers, sets, functions, relations, orders and proof strategies.
Chapter~\ref{ch:sums} develops sums, induction, binomial coefficients and finite
differences. Chapter~\ref{ch:counting} turns these into counting methods, lattice
paths, recurrences and introductory generating functions.
Chapter~\ref{ch:sequences} establishes real convergence, completeness, continuity,
series and asymptotic estimates, including the brief derivative and mean-value
toolkit needed in subsequent arguments. Chapter~\ref{ch:complex} studies complex
numbers and roots of unity. Chapter~\ref{ch:polynomials} develops polynomials,
interpolation and factorisation. Chapter~\ref{ch:linear} builds linear algebra
from spaces and maps through orthogonality, spectral theory and approximation.
Chapter~\ref{ch:calculus} gives systematic differential and integral calculus with
Taylor formulas and certified errors. Chapter~\ref{ch:inequalities} develops
inequalities and convexity as tools for estimates. Chapter~\ref{ch:probability}
uses the preceding mathematics to construct elementary probability and prove a
first law of large numbers.

The chapter order follows the dependencies between the tools. Completeness
comes before the extreme-value arguments used for polynomials, linear maps
come before determinants, and the calculus comes before convexity proved by
derivatives. A finite identity and its limiting consequence stay visibly
distinct. A few early problems use a tool proved only later; each such
problem says so. In particular, an early integral or matrix bridge can be read
for its question and completed after the relevant chapter.

\section*{The foundational contract}
A foundations book must say where its foundations begin. We use classical logic,
ordinary set notation and the standard arithmetic of the natural numbers
$\N=\{0,1,2,\ldots\}$, integers $\Z$, rationals $\Q$ and reals $\R$. Their usual
inclusions and ordered-field arithmetic are assumed. We assume well ordering of
$\N$: every nonempty subset has a least element. Chapter~\ref{ch:sums} derives the
induction principles and the elementary divisibility facts it uses.

For the reals we assume completeness: every nonempty set bounded above has a
least upper bound. Convergence criteria, intermediate and extreme values,
uniform continuity, Rolle's theorem and the mean value theorem are then proved.
We do not assume these conclusions separately. The short derivative toolkit is
in Chapter~\ref{ch:sequences}; the full rules of differentiation and their proofs
are in Chapter~\ref{ch:calculus}.

We admit the elementary functions $\exp$, $\ln$, $\cos$ and $\sin$, with their
usual real domains, addition laws and derivatives: $\exp'=\exp$, $\ln'(x)=1/x$
for $x>0$, $\sin'=\cos$ and $\cos'=-\sin$. We assume $\exp(0)=1$, that $\exp$ is
a bijection onto the positive reals with inverse $\ln$, and the usual unit-circle
normalisation and period $2\pi$ of sine and cosine. Power series used to compute
these functions are justified in the text; they are not a substitute for an
unmentioned construction of the functions themselves.

We also assume the integral of a continuous real function on a compact interval,
with linearity, order preservation, interval additivity and the integral of a
constant. The fundamental theorem, integration by parts, change of variables and
the approximation estimates used here are proved from that contract. Improper
integrals are limits, not an extra permission to interchange infinite operations.
Volume~6 develops the underlying foundations further, and Volume~4 develops the
special functions and methods of calculation in greater depth.

Probability starts with finite spaces. Countable distribution examples are marked
as an extension and use convergent sums and truncation arguments; constructing a
measure on a space of infinite sequences belongs to later volumes. Results quoted
in challenge notes as advanced outlooks are distinguished from the elementary
claims proved here and carry references. Historical sources accompany the
chapter openings; the Gauss schoolroom story, for example, is presented with its
legendary embellishments separated from the early account.

\section*{How to use the book}
Each chapter begins with a mathematical question or its history. Read the
from-scratch material with pencil in hand. In a ``Pause and think'' box, make a
prediction before uncovering the next argument. Workshops show how to discover a
proof through examples, failed approaches and a useful change of language.

Exercises have three levels: $\star$ develops technique and conceptual precision;
$\star\star$ requires a choice of method or a sequence of arguments;
$\star\star\star$ asks for a more substantial discovery. Every numbered exercise
has a separate hint and a complete solution. \emph{How to find it} explains the
search; \emph{Write-up} gives the argument that survives that search. Try the
problem before reading either. A challenge box may point to a result far beyond
the volume; its notes specify which parts are elementary and which are outlooks.

For a first reading, follow the chapter order, attempt the one-star exercises and
work through the workshops. For a second, add the two-star exercises and the
explicit bridge problems: these combine tools that initially looked unrelated.
For a preparation focused on Volume~1, give particular attention to language,
induction, counting, roots of unity, polynomial arithmetic, series and estimates;
return to the calculus and probability bridges as your interests require. The
final field guide gives an integrative investigation and a way to diagnose the
next unfamiliar problem.

\section*{From this volume to the series}
This book is Volume~1, \emph{First Principles}. The fifteen volumes form six parts.
What follows is where the tools of this volume are used again.

\begin{description}
\item[Part I. Foundations.]
\item[Volume 2: The Integers.] Proof, induction, divisibility, polynomial
arithmetic, counting and elementary asymptotic estimates.
\item[Volume 3: Symmetry.] Functions and equivalence relations, coefficient
fields, linear maps, bases, operators and roots of unity.
\item[Volume 4: Chance.] Counting, finite probability, expectation, variance,
convexity and convergence, before infinite experiments and measure.
\item[Part II. Bridges.]
\item[Volume 5: Calculation.] Series, summation by parts, Taylor errors,
integration and inequalities, used as methods rather than as foundations.
\item[Volume 6: Crossroads.] The same estimates, now at the first meeting
of arithmetic and probability.
\item[Part III. Analysis and Geometry.]
\item[Volume 7: Spaces.] Sets, functions, completeness and the integral, rebuilt
as topology, measure and limits.
\item[Volume 8: Analysis.] Limits, complex numbers and linear maps,
extended to operators, spectra and transforms.
\item[Volume 9: Shapes.] Geometry that uses the linear algebra and the estimates
of this volume as its local language.
\item[Volume 10: Abstraction.] The language of maps and equivalence, read
one level up, as categories, algebra and homology.
\item[Part IV. Randomness.]
\item[Volume 11: Randomness.] Expectation, variance and convergence, in the form
needed for martingales, limit theorems and ergodic theory.
\item[Volume 12: Random Motion.] The same probabilistic estimates along paths,
together with the spectral theorem.
\item[Part V. Number.]
\item[Volume 13: Arithmetic.] Polynomials, roots of unity and unique factorisation,
inside number fields.
\item[Volume 14: The Primes.] Series, sums and inequalities aimed at the primes.
\item[Part VI. Synthesis.]
\item[Volume 15: Order and Chance.] Finite probability and integer counting,
used to study typical integers. The subtitle is \emph{The Meeting of Structure and Randomness}.
\end{description}

The advanced theories of the later volumes are destinations, not hidden prerequisites
of this volume.
